\documentclass[11pt]{article}
%% ================================================================
%%  Preamble -- target: Nature Computational Science (Article)
%%  Limits (third-party guide, verify on nature.com):
%%    main text <= 3,500 words (excl. abstract, Methods, refs, legends)
%%    abstract <= 150 words, unreferenced
%%    <= 6 display items (figures + tables)
%% ================================================================

%% ---------- Packages ----------
\usepackage[T1]{fontenc}
\usepackage[utf8]{inputenc}
\usepackage{lmodern}
\usepackage[a4paper,margin=2.5cm]{geometry}
\usepackage{setspace}
\usepackage{amsmath, amssymb, amsthm}
\usepackage{mathtools}
\usepackage{bm}
\usepackage{graphicx}
\usepackage{subcaption}
\usepackage{booktabs}
\usepackage{xcolor}
\usepackage{algorithm}
\usepackage{algpseudocode}
\usepackage[numbers,sort&compress]{natbib}
\usepackage[colorlinks=true,linkcolor=blue!60!black,citecolor=blue!60!black,urlcolor=blue!60!black]{hyperref}
\usepackage{cleveref}   % after hyperref
\usepackage{authblk}

\graphicspath{{figures/}}

%% ---------- Notation (same as MAIN.tex and the workshop draft) ----------
\newcommand{\zstar}{z^*}
\newcommand{\ux}{u^*}
\newcommand{\uhat}{\hat{u}}
\newcommand{\bw}{\mathbf{w}}
\newcommand{\bu}{\mathbf{u}}
\newcommand{\bA}{\mathbf{A}}
\newcommand{\bb}{\mathbf{b}}
\newcommand{\R}{\mathbb{R}}
\newcommand{\Lop}{\mathcal{L}}
\newcommand{\Mset}{\mathcal{M}}          % the map (set of points)
\newcommand{\Sset}{\mathcal{S}}          % a stencil
\newcommand{\Aset}{\mathcal{A}}          % action set (menu)
\newcommand{\normone}[1]{\lVert #1 \rVert_1}
\newcommand{\normtwo}[1]{\lVert #1 \rVert_2}
\newcommand{\EPS}{\mathrm{EPS}}
\newcommand{\Nref}{N_{\mathrm{ref}}}

%% ---------- Drafting markers ----------
%% \NUM{}  = number taken from a notebook/run file; re-check before submitting
%% \CITE{} = citation not yet resolved to a bibtex key
%% \TODO{} = missing content; stays red in final mode on purpose
\newif\ifdraftmode
\draftmodetrue          % <<-- \draftmodefalse for the submission build

\newcommand{\TODO}[1]{\textcolor{red}{\textbf{[TODO: #1]}}}
\newcommand{\CITE}[1]{\textcolor{orange}{\textbf{[cite: #1]}}}
\ifdraftmode
  \newcommand{\NUM}[1]{\textcolor{blue}{\underline{#1}}}
\else
  \newcommand{\NUM}[1]{#1}
\fi

%% ---------- Theorem environments ----------
\theoremstyle{plain}
\newtheorem{proposition}{Proposition}
\theoremstyle{remark}
\newtheorem{remark}{Remark}


%% ---------- Title ----------
\title{Reinforcement Learning for Discovering Integration Schemes\\
       for Partial Differential Equations}
%% Alternative under discussion:
% \title{Learning to solve partial differential equations at a single point}

\author[1]{Ali Tayyar}
\author[1]{Joseph Bakarji\thanks{Corresponding author: \texttt{\TODO{email}}}}
\affil[1]{Department of Mechanical Engineering, American University of Beirut, Beirut, Lebanon}
\date{}

\begin{document}
\maketitle
\onehalfspacing

%% ================================================================
%%  ABSTRACT  (v1, Nov 17 2026; <= 150 words, no citations)
%% ================================================================
\begin{abstract}
Numerical methods for time-dependent partial differential equations (PDEs) compute the solution over the whole domain, even when only its value at a single point in space and time is needed. We propose a method that computes this value directly. Starting at the query point, it builds a map backward to the initial and boundary data: each point in the map chooses a stencil of points at earlier times, and only those points are ever evaluated. A reinforcement learning agent makes these choices, trading accuracy against the number of points. Its reward weighs each local error by that point's influence on the query, computed with a discrete adjoint, and never uses the exact solution. On Burgers', advection--diffusion and Kuramoto--Sivashinsky equations, the learned maps match the accuracy of classical and goal-oriented adaptive schemes with \TODO{X}-fold fewer points, on queries unseen in training. Without being told, the agent aligns its stencils with the characteristics and widens them where the equation is stiff.
\end{abstract}

%% ================================================================
%%  INTRODUCTION  (no heading in Nature format; ~600--800 words)
%%  Beats:
%%   (1) Solvers compute the whole field; many tasks need u at one point.
%%   (2) The value at one point depends only on its domain of dependence,
%%       and only part of it matters (adjoint / Green's function).
%%   (3) Goal-oriented adaptivity (DWR) exploits this, but a posteriori,
%%       on a global mesh, by iteration. Learned discretizations and
%%       RL for AMR learn resolution, but still on the whole field.
%%   (4) Here: a backward map from the query, built a priori by an agent;
%%       goal-oriented adaptivity taken to its limit.
%%   (5) Results in one paragraph; the physics the agent recovers.
%% ================================================================
\TODO{introduction}

%% ================================================================
%%  RESULTS  (one subsection per main figure; <= 6 display items)
%% ================================================================
\section*{Results}

\subsection*{Computing the solution at one point}
\label{sec:res-map}
%% Fig. 1: query -> backward map -> influence beta; error identity.
\TODO{backward map, influence, error identity}

\subsection*{An agent that builds the map}
\label{sec:res-agent}
%% Fig. 2: two passes (coarse lattice, then agent); MDP in one schematic.
\TODO{two-pass method and the decision process}

\subsection*{Fewer points at equal accuracy}
\label{sec:res-cost}
%% Fig. 3: error at z* vs points, three PDEs, against baselines.
\TODO{cost--accuracy results}

\subsection*{The learned stencils follow the characteristics}
\label{sec:res-char}
%% Fig. 4: beta and stencil shifts vs characteristic / adjoint; null control.
\TODO{physical evidence I}

\subsection*{Stiffness sets the stencil width}
\label{sec:res-stiff}
%% Fig. 5: KS, wide spacing, instability of compact stencils.
\TODO{physical evidence II}

\subsection*{Generalization to unseen queries}
\label{sec:res-gen}
%% Fig. 6: interpolation vs extrapolation in t*.
\TODO{held-out queries}

%% ================================================================
%%  DISCUSSION  (meaning, limitations, outlook)
%% ================================================================
\section*{Discussion}
\TODO{discussion, limitations, outlook}

%% ================================================================
%%  METHODS  (after the main text in Nature format; full detail)
%% ================================================================
\section*{Methods}

\subsection*{Problem setting}
\label{sec:m-setting}
\TODO{PDE class, grid, query $\zstar=(x^*,t^*)$, known data}

\subsection*{Taylor weights and admissible stencils}
\label{sec:m-weights}
\TODO{min-norm weights, rows, admissibility}

\subsection*{The backward map and the forward pass}
\label{sec:m-map}
\TODO{expansion order, forward evaluation, Picard re-solve}

\subsection*{Influence and the error identity}
\label{sec:m-beta}
\TODO{$\beta$ recursion, $e(\zstar)=\sum_p \beta_p\delta_p$, bound}

\subsection*{A-priori error score}
\label{sec:m-score}
\TODO{Taylor residual with guessed derivatives}

\subsection*{Coarse first pass}
\label{sec:m-pass1}
\TODO{uniform backward lattice, local polynomial fit}

\subsection*{Stencil menu}
\label{sec:m-menu}
\TODO{admissibility and stability rule}

\subsection*{Decision process and training}
\label{sec:m-mdp}
\TODO{state, action, reward, MaskablePPO, multi-query training}

\subsection*{Test problems}
\label{sec:m-pdes}
\TODO{Burgers, advection--diffusion, Kuramoto--Sivashinsky}

\subsection*{Baselines and metrics}
\label{sec:m-baselines}
\TODO{FD-3, LW, DWR, ETDRK4; equal-cost comparison}

%% ================================================================
%%  END MATTER
%% ================================================================
\section*{Data availability}
\TODO{data statement}

\section*{Code availability}
\TODO{repository link}

\bibliographystyle{unsrtnat}
\bibliography{reference}

\section*{Acknowledgements}
\TODO{funding}

\section*{Author contributions}
\TODO{CRediT statement}

\section*{Competing interests}
The authors declare no competing interests.

\end{document}
